% !TEX encoding = UTF-8
% !TEX program = lualatex

\chapter*{СЛОВАРЬ ТЕРМИНОВ}
\addcontentsline{toc}{chapter}{СЛОВАРЬ ТЕРМИНОВ}
\markboth{СЛОВАРЬ ТЕРМИНОВ}{СЛОВАРЬ ТЕРМИНОВ}

\begin{description}
    \item[Атрибутивный метаграф] — расширение метаграфа, в котором вершинам и гипердугам сопоставлены атрибуты, содержащие семантическую, временную или количественную информацию.
    \item[Метаграф] — обобщение гиперграфа, в котором как источники, так и цели дуги могут быть множествами вершин.
    \item[Гипердуга] — дуга метаграфа, соединяющая множество источников с множеством целей.
    \item[Цепочка компрометации] — последовательность действий злоумышленника, направленных на достижение цели атаки (например, кража данных).
    \item[TTP (Tactics, Techniques, Procedures)] — тактики, техники и процедуры, используемые злоумышленниками; стандартизированы в фреймворке MITRE ATT\&CK.
    \item[MITRE ATT\&CK] — глобальная база знаний о поведении киберугроз, структурированная по тактикам и техникам.
    \item[APT (Advanced Persistent Threa)] — целевая многоэтапная атака, осуществляемая высококвалифицированными акторами с длительным присутствием в инфраструктуре жертвы.
    \item[EDR (Endpoint Detection and Response)] — система обнаружения и реагирования на угрозы на конечных устройствах.
    \item[SIEM (Security Information and Event Management)] — платформа для сбора, анализа и корреляции событий безопасности.
    \item[SOAR (Security Orchestration, Automation and Response)] — платформа автоматизации процессов реагирования на инциденты.
    \item[Телеметрия безопасности] — данные о событиях безопасности, поступающие от источников (EDR, сетевые сенсоры, аудит ОС).
    \item[Нормализация событий] — приведение событий из разнородных источников к единой схеме (например, UCOS).
    \item[Динамический метаграф] — метаграф, обновляемый инкрементально по мере поступления новых событий телеметрии.
    \item[Эталонный шаблон атаки] — атрибутивный метаграф, описывающий известную цепочку компрометации (например, APT29).
    \item[Частичный изоморфизм метаграфов] — соответствие между подмножеством вершин и дуг эталонного и динамического метаграфов с учётом атрибутов.
    \item[Мера риска компрометации] — количественная оценка угрозы, основанная на вероятности реализации критических путей в метаграфе.
    \item[Латеральное перемещение] — тактика перемещения злоумышленника по сети от одного скомпрометированного хоста к другому.
    \item[Персистентность] — тактика обеспечения устойчивого присутствия в системе после первоначальной компрометации.
    \item[C2 (Command and Control)] — канал связи между скомпрометированной системой и инфраструктурой злоумышленника.
    \item[Living-off-the-land] — техники использования легитимных системных инструментов (PowerShell, WMI) для маскировки атаки.
    \item[STIX (Structured Threat Information eXpression)] — стандарт обмена информацией об угрозах в машиночитаемом формате.
    \item[TAXII (Trusted Automated eXchange of Indicator Information)] — протокол передачи данных STIX.
    \item[UCOS (Unified Cyber Observation Schema)] — унифицированная схема представления событий безопасности.
    \item[Граф атак] — ориентированный граф, вершины которого представляют состояния системы, а рёбра — действия злоумышленника.
    \item[Дерево атак] — иерархическая модель атаки, где корень — цель, а листья — примитивные действия.
    \item[Многоагентная система] — система, состоящая из автономных агентов, взаимодействующих для достижения общей цели (в контексте ИБ — моделирование защиты или атаки).
\end{description}